\documentclass[11pt]{article}
\usepackage[a4paper,margin=25mm]{geometry}
\usepackage{amsmath,amssymb,amsthm,microtype}
\usepackage[hidelinks]{hyperref}
\newtheorem{theorem}{Theorem}\newtheorem{lemma}[theorem]{Lemma}
\newcommand{\E}{\mathbb E}\newcommand{\Pp}{\mathbb P}
\title{The right-strip comparison from five conditional counts}
\author{}\date{28 October 2026}
\begin{document}\maketitle

We use the five-count reserve of \texttt{loss-direct.tex} to bound the full
signed Berry--Esseen excess. Interpolating two binomial CDF values is
less expensive than comparing to one prescribed integer anchor. In
particular the normal error is a signed interpolation remainder of order
$1/n$. No geometric budget or displacement bound of order $1/n$ is needed.

\begin{theorem}\label{thm:right}
Let $n\ge200$ be an integer and $p\in[2/5,53/100]$, with $q=1-p$ and
$v=pq$. Let $Y=B+U$, where $B$ is Bernoulli of parameter $p$,
$\E(U\mid B)=0$, and the conditional noise laws satisfy
\[
 -3/5\le U_0\le a_0,\qquad -a_1\le U_1\le3/5,
 \qquad a_0,a_1\ge0,\qquad a_0+a_1\le9/125.
\]
Put $s=\E U^2$, $\lambda=ns$, $X=(Y-p)/\sqrt{v+s}$ and
$\rho=\E|X|^3$. If $3/40\le\lambda\le9/25$, then for independent
copies $X_1,\ldots,X_n$ and every $|H|\le46/25$,
\[
 \sqrt n\{\Pp(X_1+\cdots+X_n\le H)-\Phi(H/\sqrt n)\}
                 -C_E\rho<-\frac{\lambda}{100},
 \qquad C_E=\frac{3+\sqrt{10}}{6\sqrt{2\pi}}.
\]
\end{theorem}

This is a theorem about every law satisfying the displayed support,
centering, variance and threshold hypotheses. Acquiring these hypotheses
for an extremal witness is a separate task. The statement concerns the
positive signed discrepancy; reflection requires the corresponding
hypotheses for the reflected law. It is not an unconditional cutoff
theorem for the full iid conjecture.

\section{The crossing formula and five binomial masses}
Write
\[
 t=np+H\sqrt{v+s},\quad k=\lfloor t\rfloor,\quad u=t-k\in[0,1),
 \quad w=t-np,
\]
and let $\pi_j$ and $F_B$ be the mass function and CDF of
$\operatorname{Bin}(n,p)$. The support hypotheses imply the interior
crossing conditions $na_0\le t<n(1-a_1)$.
Indeed $v\ge6/25$, $\sqrt{v+s}<251/500$, and
\begin{equation}\label{eq:w}
 |w|<231/250<47/50.
\end{equation}
The retained v13 crossing identity gives
\begin{equation}\label{eq:cross}
 \Pp\Big(\sum X_i\le H\Big)
 =(1-u)F_B(k)+uF_B(k+1)+\mathcal C(t)-\sum_{i=1}^n\pi_i g_i(t),
\end{equation}
where $g_i\ge0$ are the kernels in \texttt{loss-direct.tex} and
\[
 \mathcal C(t)=\sum_{j=1}^{n-1}(\pi_{j+1}-\pi_j)
 \left\{\E(t-j-W_j)_+-(t-j)_+\right\}.
\]
Here $W_j$ is the sum of $j$ independent $U_1$ variables and $n-j$
independent $U_0$ variables. All lower threshold events are closed.

For $k-1\le i\le k+3$, \eqref{eq:w} gives $|i-np|<4$.
For $z=i/n$, every point on the segment joining $z$ to $p$ lies in
$[.38,.55]$, where $z(1-z)>23/100$. In particular the counts are interior.
The explicit finite Stirling remainder
\[
 \log(m!)=(m+\tfrac12)\log m-m+\tfrac12\log(2\pi)+r_m,
 \qquad0<r_m<1/(12m),
\]
and $D''(z\Vert p)=1/(z(1-z))$ give
\begin{equation}\label{eq:mass}
 \sqrt n\pi_i\ge\frac{\phi(0)}{\sqrt{z(1-z)}}
 \exp\left\{-\frac{8+1/12}{200(23/100)}\right\}
 >\frac{199}{250}\left(1-\frac{8+1/12}{46}\right)>
 \frac{13}{20}.
\end{equation}
The same finite Stirling proof as in the retained count appendix is used;
no local central limit theorem is invoked. The elementary bounds
$199/500<\phi(0)<2/5$ follow from $25/8<\pi<22/7$.

Let $\alpha=39/40$, $b=203/200$ and $x=3\alpha\lambda$.
Taylor's integral remainder gives
$1-e^{-x}\ge x-x^2/2+x^3/6-x^4/24$ for every $x\ge0$.
The five-count theorem, monotonicity in $n$, and \eqref{eq:mass} yield
\begin{equation}\label{eq:W}
 \sqrt n\sum_i\pi_i g_i\ge W(\lambda):=
 \frac{13}{50}\left\{
 \frac{x-x^2/2+x^3/6-x^4/24}{2}
 -\frac{(9/25)b\lambda+3b^2\lambda^2}{16}
 -\frac{10(22/7)^2}{200}\lambda\right\}.
\end{equation}
No omitted count needs a new estimate: all other weighted kernels are
nonnegative. The finite-window theorem already controls its omitted bump
translates, with their entire fourth-moment tail.

\section{Curvature and the signed normal interpolation}
Use the summable-curvature lemma in the v14 note
\texttt{geometry-slack.tex}. Set
\[
 c=(np-q-t)_+,\qquad T=\frac{100}{47}\lambda,
 \qquad M=\frac9{25}T+3T^2.
\]
The actual conditional variance on the increasing side of the binomial
mass is at most $\lambda/q\le T$, and its fourth moment is at most $M$.
For $H\le0$ we have $c<23/50$; for $H\ge0$ it is zero.
Also
\[
 \sqrt T<\frac32\sqrt\lambda\le\frac32(\lambda+1/4),
 \qquad v^{3/2}>117/1000.
\]
The last square-root inequality is $(\sqrt\lambda-1/2)^2\ge0$.

The binomial mode coefficient $\beta=41/100$ in the curvature lemma
is valid throughout this right range. To check it, put $N=n+1$ and
$r=m/N$ for a mode $m=\lfloor(N+1)p\rfloor$. Since
$|m-Np|\le.6$ and $|1-2p|\le.2$,
\[
 r(1-r)\ge v(1-1/N-2/N^2)>39v/40\qquad(N\ge43).
\]
Finite Stirling, discarding its nonpositive entropy and denominator
remainders, gives
\[
 \max_j\Pp(\operatorname{Bin}(N,p)=j)
 <\frac{(2/5)\sqrt{40/39}(516/515)}{\sqrt{Nv}}
 <\frac{41/100}{\sqrt{Nv}}.
\]
Thus, writing
\[
 S(\lambda,c)=\frac T4+\frac{3c}{2}(\lambda+1/4)
       +\frac{27M}{256}\left(\frac53+\frac{5c}{4}\right),
\]
we have $\sqrt n\mathcal C(t)\le(41/100)S(\lambda,c)/(nv^{3/2})$.

Schulz's binomial theorem \cite[Theorem1, printed p.1]{Schulz} gives,
at both integer anchors and for every $0<p<1$,
\[
 \sqrt n\{F_B(j)-\Phi((j-np)/\sqrt{nv})\}<C_E\rho_B,
 \qquad \rho_B=(1-2v)/\sqrt v.
\]
Convex interpolation preserves this bound. No sharper middle-$p$ reserve
is needed here. For $G(z)=\Phi((z-np)/\sqrt{nv})$,
\[
 G''(z)\le\frac{\phi(0)(np-z)_+}{(nv)^{3/2}}.
\]
The signed linear-interpolation remainder therefore satisfies
\[
 \sqrt n\{(1-u)G(k)+uG(k+1)-G(t)\}
 \le\frac{\phi(0)}{2nv^{3/2}}u(1-u)(u-w)_+.
\]
For $H\le0$, the variance change to $v+s$ is favorable. Using
$u(1-u)\le1/4$, $u^2(1-u)\le4/27$ and \eqref{eq:w}, the total
normal cost is at most
\[
 \frac{2/5}{nv^{3/2}}\left(\frac{231/250}{8}+\frac2{27}\right).
\]
For $H\ge0$, the interpolation cost is at most
$(2/5)(2/27)/(nv^{3/2})$ and the variance change costs at most
$\phi(0)ws/(2v^{3/2})$. The total is at most
\[
 \frac{2/5}{nv^{3/2}}\left(\frac2{27}+\frac{231}{500}\lambda\right).
\]
These estimates retain both signs of $H$, every phase and all threshold atoms.

The cubic correction has the favorable sign. Expanding the actual marginal
absolute third moment, and using the inward cubic bounds, gives
\[
 (v+s)^{3/2}(\rho-\rho_B)
 \ge s\left\{\frac{87}{250}-\frac{13}{16}s\right\}>0.
\]
Here the two conditional supports have fixed signs relative to $p$;
$\E U_0^3\le a_0\E U_0^2$, $\E U_1^3\ge-a_1\E U_1^2$,
and $(1+y)^{3/2}\le1+3y/2+3y^2/8$ for $y\ge0$.
The right side is positive because $s\le9/5000$.

\section{Four rational endpoint comparisons complete the proof}
Define two quadratic error polynomials
\begin{align*}
 E_-(\lambda)&=\frac{(41/100)S(\lambda,23/50)
 +(2/5)(231/2000+2/27)}{117/1000},\\
 E_+(\lambda)&=\frac{(41/100)S(\lambda,0)
 +(2/5)(2/27+231\lambda/500)}{117/1000}.
\end{align*}
The preceding estimates give the following strict upper bound on the full signed excess:
$-W(\lambda)+E_\pm(\lambda)/n$, according to the sign of $H$.
It remains to prove positivity of
\[
 P_\pm(\lambda)=W(\lambda)-E_\pm(\lambda)/200-\lambda/100.
\]
Both are quartics. Their concavity has a direct explanation. With
$z=3\alpha$, differentiating the gain gives
\[
 W''(\lambda)=-\frac{13}{50}
 \left\{\frac{z^2}{4}\bigl[(z\lambda-1)^2+1\bigr]
                    +\frac{3b^2}{8}\right\}<0.
\]
The error polynomials $E_\pm$ are convex quadratics. Thus $P_\pm''<0$
for every real $\lambda$.
The minimum of each concave polynomial is attained at an endpoint.
The four exact rational comparisons, rounded down only for display, are
\[
\begin{array}{c|cc}
 &\lambda=3/40&\lambda=9/25\\\hline
P_-&>6075/10^6&>2577/10^6\\
P_+&>11528/10^6&>11247/10^6.
\end{array}
\]
The companion checker reconstructs the polynomials and verifies these
comparisons using fractions. It performs no subdivision or optimization.
Thus $P_\pm>0$ on the entire closed noise interval. Since $E_\pm>0$
and $n\ge200$, Theorem~\ref{thm:right} follows.

\section*{Dependencies and remaining range}
The new local inputs are the five-count reserve in \texttt{loss-direct.tex},
the v13 crossing identity and the v14 summable-curvature lemma. The external
probability input is precisely Schulz's all-$p$ binomial theorem; finite
Stirling has its explicit retained elementary proof. No geometric-budget,
selected-count transport, marked crossing, Poisson comparison or new
parameter certificate is used in this larger-noise comparison.

For acquired witnesses, v13 geometry supplies $\lambda<.36$, while v12
handles smaller positive noise. That small-noise theorem may choose a
different integer anchor; this causes no problem because the binomial
theorem holds at every anchor. Zero noise is Bernoulli. The support strips,
raw-threshold bound, sign-specific acquisition and independent finite base
remain separate obligations. No statement about unrestricted $C_n$ is made.

\begin{thebibliography}{1}
\bibitem{Schulz}J. Schulz,
\emph{The optimal Berry--Esseen constant in the binomial case},
doctoral dissertation, Universit\"at Trier, 2016.
\url{https://d-nb.info/1197702695/34}.
\end{thebibliography}
\end{document}

% Deliberate reviewed-source mutation.
